% ============================================================
% portada.tex — Portadilla textual del volumen
% Imita la composición tipográfica de la portada original de 1889
% (no la pasta física del libro), en español y paleta azul/rojo.
% ============================================================
\begin{titlepage}
\centering
\vspace*{1.5cm}

{\color{azulpoincare}\small\scshape Curso de la Facultad de Ciencias de París\par}
{\color{grisnota}\footnotesize publicado por la asociación amical de alumnos y exalumnos\par
de la Facultad de Ciencias\par}

\vspace{2pt}
\color{rojopoincare}\rule{0.55\textwidth}{0.6pt}
\color{black}
\vspace{18pt}

{\color{azulpoincare}\scshape\large Curso de Física Matemática\par}

\vspace{28pt}

{\color{azulpoincare}\scshape\Large Teoría Matemática\par}
{\color{azulpoincare}\scshape\large de\par}
{\color{azulpoincare}\bfseries\scshape\Huge La Luz\par}

\vspace{10pt}
{\color{rojopoincare}\large I\par}

\vspace{20pt}

{\small Volumen I — Edición guía en español latino\par}
{\footnotesize Traducción, desarrollo explícito de cálculos y figuras vectoriales\par
a partir de las lecciones profesadas por H.\,Poincaré durante\par
el primer semestre de 1887-1888\par}

\vspace{40pt}

{\scshape por\par}
\vspace{4pt}
{\Large\bfseries H.\ Poincaré}, {\small miembro del Instituto\par}

\vspace{6pt}
{\footnotesize redactadas originalmente por J.\ Blondin, licenciado en ciencias\par}

\vfill

\color{rojopoincare}\rule{0.35\textwidth}{0.6pt}\color{black}

\vspace{14pt}
{\footnotesize Edición y traducción al español latino, \the\year\par}

\end{titlepage}
